\chapter{Analiza problemei}
\etapa{S1-E3}

\section{Descrierea activităților desfășurate}
În cadrul acestei etape, activitatea s-a desfășurat pe parcursul a două săptămâni:
\begin{itemize}
    \item \textbf{Săptămâna 1 -- Analiza documentației și a cerințelor:} S-a efectuat studiul documentului introductiv de testare software pentru înțelegerea conceptelor fundamentale: ciclul de viață al dezvoltării software (SDLC), metodologii Agile/Scrum/Kanban, principii de testare, niveluri și tipuri de testare (Black-box, Functional, Non-functional, Regresie, Smoke testing), precum și structura unui defect (bug) și a cazurilor de testare (Test Cases).
    \item \textbf{Săptămâna 2 -- Testare dinamică Exploratorie și Funcțională:} S-a primit accesul la aplicația/site-ul web alocat pentru testare. S-a realizat o sesiune completă de testare exploratorie și black-box, interacționând cu toate elementele interfeței (butoane, formulare, meniuri, linkuri) pentru a verifica conformitatea comportamentului aplicației cu cerințele funcționale. S-au identificat, izolat și raportat bug-urile și anomaliile descoperite, respectând structura standard a unui tichet de tip bug (Titlu, Pași de reproducere, Comportament observat vs. Comportament așteptat, Severitate/Prioritate și atașamente doveditoare).
\end{itemize}

\section{Cerințele soluției}
\begin{enumerate}
    \item Înțelegerea și însușirea conceptelor teoretice și a terminologiei specifice asigurării calității (QA) din documentația de referință.
    \item Aplicarea practică a tehnicilor de testare manuală (Black-box și Exploratorie) pe aplicația indicată.
    \item Documentarea corectă a defectelor găsite, utilizând componentele standard ale unui bug (pași clari de reproducere, rezultat actual vs. rezultat așteptat).
\end{enumerate}

\section{Datele utilizate}
\begin{itemize}
    \item \textbf{Date de intrare:} Documentația tehnică de testare și link-ul/accesul la aplicația web pusă la dispoziție pentru testare.
    \item \textbf{Date de test:} Date valide și invalide introduse în câmpurile aplicației (formulare, butoane de autentificare/navigare) pentru a verifica comportamentul sistemului în diferite scenarii.
\end{itemize}

\section{Constrângerile proiectului}
\begin{itemize}
    \item Testarea s-a limitat la interfața de utilizator (UI) și la comportamentul extern al aplicației (Black-box testing), fără a avea acces direct la modificarea codului sursă.
    \item Timpul alocat a fost fixat pe parcursul celor două săptămâni ale etapei (o săptămână pentru analiză/teorie și o săptămână pentru execuție și identificare de bug-uri).
\end{itemize}

\section{Criteriile de evaluare a soluției}
\begin{itemize}
    \item \textbf{Corectitudinea identificării defectelor:} Capacitatea de a semnala corect erorile reale din aplicație și de a diferenția problemele majore de cele cosmetice (Trivial/Minor vs. Critical/Blocker).
    \item \textbf{Calitatea raportării:} Claritatea pașilor de reproducere a bug-urilor și acuratețea descrierii rezultatului observat comparativ cu cel așteptat.
    \item \textbf{Acoperirea testării:} Gradul în care s-a parcurs întreaga aplicație (toate butoanele, ecranele și funcționalitățile principale vizibile).
\end{itemize}